\section{TensorFlow Model Implementation}

This section describes the TensorFlow/Keras implementation of the keyword detection model, which was collaborated on by the secondary project developer. This model differs primarily from the PyTorch model by integrating a custom Mel-Spectrogram layer directly into the neural network graph, allowing it to process raw audio waveforms on the GPU.

\subsection{Dataset Pipeline and Preprocessing}
The TensorFlow pipeline uses native TensorFlow operations (\texttt{tf.data} and \texttt{tf.io}) to load and decode audio files. 

\begin{lstlisting}[language=Python, caption={Native TensorFlow WAV Decoding and Alignment}, label={lst:tf_process_path}]
@tf.function
def process_path(file_path, label):
    # Read and decode raw WAV binary
    audio_binary = tf.io.read_file(file_path)
    audio, _ = tf.audio.decode_wav(audio_binary, desired_channels=1)
    audio = tf.squeeze(audio, axis=-1)

    # Pad or crop to exactly 16000 samples (1 second)
    audio_length = tf.shape(audio)[0]
    padding = tf.maximum(NUM_SAMPLES - audio_length, 0)
    audio = tf.pad(audio, [[0, padding]])
    audio = audio[:NUM_SAMPLES]

    return audio, label
\end{lstlisting}

\subsubsection{Dataset Splits}
Unlike the PyTorch model's oversampling and downsampling balancing steps, the TensorFlow pipeline uses a standard stratified 3-way split of the raw dataset:
\begin{itemize}
    \item \textbf{Training Set:} 70\% of the files.
    \item \textbf{Validation Set:} 10\% of the files.
    \item \textbf{Test Set:} 20\% of the files.
\end{itemize}

\subsection{Data Augmentation}
Audio augmentation is applied dynamically to the raw audio waveforms on the GPU using TensorFlow tensors:
1. **Random Gaussian Noise:** Injects a standard normal distribution of noise to the waveform:
   $$\tilde{x} = x + \mathcal{N}(0, 0.003)$$
2. **Random Time Shift:** Randomly shifts the waveform array using circular shift:
   $$\tilde{x} = \text{roll}(x, \text{shift}) \quad \text{where } \text{shift} \in [-1000, 1000] \text{ samples}$$
3. **Volume Variation (Gain):** Scales the amplitude by a random factor:
   $$\tilde{x} = x \times g \quad \text{where } g \sim \mathcal{U}(0.8, 1.2)$$

\subsection{Onboard Mel-Spectrogram Layer}
The model includes a custom Keras layer, \texttt{MelSpectrogram}, which computes the Log-Mel Spectrogram of the raw input waveform. This design choice allows the model to accept raw audio float32 arrays directly, rather than requiring separate preprocessing transforms during deployment.

The transformation parameters are configured to match a 64 mel filterbank:
* **Frame Length:** $1024$ samples ($64\text{ ms}$)
* **Frame Step:** $512$ samples ($32\text{ ms}$)
* **FFT Length:** $1024$
* **Mel Bins:** $64$
* **Frequency Range:** $80\text{ Hz}$ to $7600\text{ Hz}$

\begin{lstlisting}[language=Python, caption={Custom Keras Onboard MelSpectrogram Layer}]
class MelSpectrogram(tf.keras.layers.Layer):
    def __init__(self, **kwargs):
        super(MelSpectrogram, self).__init__(**kwargs)

    def call(self, audio):
        # Compute Short-Time Fourier Transform (STFT)
        stfts = tf.signal.stft(audio, frame_length=1024,
                               frame_step=512, fft_length=1024)
        spectrograms = tf.abs(stfts)
        
        # Warp the linear scale spectrograms into the mel-scale
        num_spectrogram_bins = stfts.shape[-1]
        mel_weight_matrix = tf.signal.linear_to_mel_weight_matrix(
            num_mel_bins, num_spectrogram_bins, TARGET_SAMPLE_RATE,
            lower_edge_hertz, upper_edge_hertz)
        
        mel_spectrograms = tf.tensordot(spectrograms, mel_weight_matrix, 1)
        mel_spectrograms.set_shape(spectrograms.shape[:-1].concatenate(
            mel_weight_matrix.shape[-1:]))
            
        # Log mel spectrograms
        log_mel_spectrograms = tf.math.log(mel_spectrograms + 1e-6)
        
        # Add channel dimension: (Batch, Frames, Mels, 1)
        return tf.expand_dims(log_mel_spectrograms, -1)
\end{lstlisting}

\subsection{Model Architecture}
The Keras model takes an input shape of $(Batch, 16000)$ and feeds it through the custom MelSpectrogram layer before passing it to the 2D Convolutional layers.

\begin{lstlisting}[language=Python, caption={TensorFlow/Keras CNN Model Structure}]
def build_model(num_classes):
    audio_input = tf.keras.Input(shape=(NUM_SAMPLES,), name='audio')
    
    # Extract Mel-Spectrogram on-board
    x = MelSpectrogram()(audio_input)

    # Convolutional Layers
    x = tf.keras.layers.Conv2D(16, 3, padding='same', activation='relu')(x)
    x = tf.keras.layers.MaxPooling2D(2, 2)(x)
    
    x = tf.keras.layers.Conv2D(32, 3, padding='same', activation='relu')(x)
    x = tf.keras.layers.MaxPooling2D(2, 2)(x)
    
    x = tf.keras.layers.Conv2D(64, 3, padding='same', activation='relu')(x)
    x = tf.keras.layers.MaxPooling2D(2, 2)(x)
    
    # Resize spatial features to a fixed 4x4 resolution
    x = tf.keras.layers.Resizing(4, 4)(x)
    
    # Dense Fully Connected Layers
    x = tf.keras.layers.Flatten()(x)
    x = tf.keras.layers.Dense(128, activation='relu')(x)
    x = tf.keras.layers.Dropout(0.3)(x)
    outputs = tf.keras.layers.Dense(num_classes)(x)
    
    return tf.keras.Model(inputs=audio_input, outputs=outputs)
\end{lstlisting}

Unlike the PyTorch model which uses adaptive average pooling to enforce output dimensions, this network uses the Keras \texttt{Resizing} layer to resize spatial features to a fixed $4\times4$ resolution, providing a deterministic input size of 1024 flat units before the dense classification layer.

\subsection{Training Configurations}
The model is compiled with the Adam optimizer ($\eta = 0.001$) and trained over 30 epochs:
\begin{lstlisting}[language=Python, caption={Compile and Callbacks Configuration}]
model.compile(
    optimizer=tf.keras.optimizers.Adam(learning_rate=0.001),
    loss='sparse_categorical_crossentropy',
    metrics=['accuracy']
)

early_stopping = tf.keras.callbacks.EarlyStopping(
    monitor='val_loss',
    patience=5,
    restore_best_weights=True
)

reduce_lr = tf.keras.callbacks.ReduceLROnPlateau(
    monitor='val_loss',
    factor=0.5,
    patience=2,
    verbose=1
)
\end{lstlisting}
*Note: Using \texttt{'sparse\_categorical\_crossentropy'} on a linear output without a softmax activation layer requires setting \texttt{from\_logits=True} during compilation to prevent training instabilities. Although the notebook omitted this parameter, it serves as a point of comparison with PyTorch's cross-entropy loss.*

\subsection{ONNX Export via tf2onnx}
After training, the model is evaluated on the test set and converted to ONNX using the \texttt{tf2onnx} package:
\begin{lstlisting}[language=Python, caption={tf2onnx Model Conversion}]
import tf2onnx

# Define dynamic batch input signature
input_signature = [tf.TensorSpec([None, NUM_SAMPLES], tf.float32, name='input')]

# Convert from Keras model
onnx_model, _ = tf2onnx.convert.from_keras(
    model,
    input_signature=input_signature,
    opset=13
)

with open('../keyword_detector.onnx', "wb") as f:
    f.write(onnx_model.SerializeToString())
\end{lstlisting}
This packages the custom MelSpectrogram layer directly into the exported graph, saving the model as a self-contained \texttt{keyword\_detector.onnx} file that accepts raw float arrays.
